\SourceLeaf{219}%
\noindent If we now hold \(x\) and \(y\) fixed and vary \(l,\lambda\), differentiation gives
\[
\begin{aligned}
\frac{\mathrm{d}l}{\mathrm{d}\lambda}&=\frac{ab'-ba'}{(a\cos\lambda+a'\sin\lambda)^2+(b\cos\lambda+b'\sin\lambda)^2}\\
&=(ab'-ba')\mathord{.}\frac{\mathrm{d}\sigma^2}{\mathrm{d}s^2}
\end{aligned}
\]
\noindent We see that according as \(ab'-ba'\) is positive \textit{or} negative, \(l\) and \(\lambda\) always increase together \textit{or} change in opposite directions. In the first case representations 2 and 3 are similarly positioned, and in the other their positions are reversed.

Combining this result with the preceding one, we find that representations 1 and 3 are similarly positioned or reversed according as \(\frac{ab'-ba'}{h}\) is positive or negative.

On the surface with equation \(\psi=0\) we have
\[
e\,\mathrm{d}x+g\,\mathrm{d}y+h\,\mathrm{d}z=0
\]
\noindent and hence also
\[
(ea+gb+hc)\mathrm{d}t+(ea'+gb'+hc')\mathrm{d}u=0
\]
\noindent whatever ratio we choose for \(\mathrm{d}t\) and \(\mathrm{d}u\). Thus the following identities must hold
\[
ea+gb+hc=0,\qquad ea'+gb'+hc'=0
\]
\noindent from which it follows that \(e,g,h\) are proportional to the quantities \(bc'-cb',ca'-ac',ab'-ba'\), respectively. Hence
\[
\frac{bc'-cb'}{e}=\frac{ca'-ac'}{g}=\frac{ab'-ba'}{h}
\]
\noindent We can therefore take any one of these three expressions or multiply by the inherently positive quantity \(ee+gg+hh\) and take the resulting symmetric quantity
\[
ebc'+gca'+hab'-ecb'-gac'-hba'
\]
\noindent as the criterion of similar or reversed position of the parts in representations 1 and 3.

In exactly the same way, the similar or reversed position of the parts in representations 6 and 8 depends on whether the quantity
\[
\frac{BC'-CB'}{E}=\frac{CA'-AC'}{G}=\frac{AB'-BA'}{H}
\]
\SourceLeaf{220}%
\noindent or, if preferred, the symmetric quantity
\[
EBC'+GCA'+HAB'-ECB'-GAC'-HBA'
\]
\noindent is positive or negative.

The comparison of representations 3 and 4 proceeds on much the same grounds as that of 2 and 3. Whether the parts are similarly positioned or reversed depends on the positive or negative sign of the quantity
\[
(\frac{\mathrm{d}p}{\mathrm{d}t})\mathord{.}(\frac{\mathrm{d}q}{\mathrm{d}u})-(\frac{\mathrm{d}p}{\mathrm{d}u})\mathord{.}(\frac{\mathrm{d}q}{\mathrm{d}t})
\]
\noindent and in the same way the positive or negative sign of
\[
(\frac{\mathrm{d}P}{\mathrm{d}T})\mathord{.}(\frac{\mathrm{d}Q}{\mathrm{d}U})-(\frac{\mathrm{d}P}{\mathrm{d}U})\mathord{.}(\frac{\mathrm{d}Q}{\mathrm{d}T})
\]
\noindent determines the similar or reversed position of the parts in representations 5 and 6.

Finally, to compare representations 4 and 5 we can refer to the analysis in Article 8. It shows that their smallest parts have similar or reversed positions according as we choose the first or second solution, that is, according as we set
\[
P+iQ=f(p+iq)\quad\text{and}\quad P-iQ=f'(p-iq)
\]
\noindent or
\[
P+iQ=f(p-iq)\quad\text{and}\quad P-iQ=f'(p+iq)
\]
\noindent respectively.

From all this we conclude that, for the smallest parts of the representation on the surface with equation \(\Psi=0\) to be similar in shape and also similarly positioned to the original on the surface with equation \(\psi=0\), we must count the negative quantities among these four
\[
\begin{gathered}
\frac{ab'-ba'}{h},\quad (\frac{\mathrm{d}p}{\mathrm{d}t})\mathord{.}(\frac{\mathrm{d}q}{\mathrm{d}u})-(\frac{\mathrm{d}p}{\mathrm{d}u})\mathord{.}(\frac{\mathrm{d}q}{\mathrm{d}t}),\\
(\frac{\mathrm{d}P}{\mathrm{d}T})\mathord{.}(\frac{\mathrm{d}Q}{\mathrm{d}U})-(\frac{\mathrm{d}P}{\mathrm{d}U})\mathord{.}(\frac{\mathrm{d}Q}{\mathrm{d}T}),\quad \frac{AB'-BA'}{H}
\end{gathered}
\]
\noindent If there are none or an even number, we must choose the first solution. If one or three of them are negative, we must choose the second. The opposite choice always produces a reversed similarity.
\par\indent\SourceLeaf{221}%
It can also be shown that if we denote the four quantities above by \(r,s,S,R\), respectively, then we always have
\[
\frac{r\BareRoot{(ee+gg+hh)}}{s}=\pm n,\qquad\frac{R\BareRoot{(EE+GG+HH)}}{S}=\pm N
\]
\noindent with \(n\) and \(N\) in the sense of Article 5. We omit the proof of this theorem here, although it is easy to find, since we have no further need of it for our purpose.

